% !TEX program = xelatex
% Example: Basic Beamer document with default theme
% Demonstrates a title page, an outline frame, and content frames
% with section divisions and \frametitle.
\documentclass{beamer}
\usetheme{default}

\title{Introduction to Beamer}
\subtitle{Creating Slides with LaTeX}
\author{Dr.~Alan Turing}
\institute{University of Computing}
\date{2025}

\begin{document}

\begin{frame}
  \titlepage
\end{frame}

\begin{frame}{Outline}
  \tableofcontents
\end{frame}

\section{Motivation}

\begin{frame}{Why Beamer?}
  \begin{itemize}
    \item Professional, print-quality slides.
    \item Seamless integration with math and TikZ.
    \item Version-controllable plain-text source.
  \end{itemize}
\end{frame}

\section{Workflow}

\begin{frame}{A Typical Workflow}
  \begin{enumerate}
    \item Write the \texttt{.tex} source.
    \item Compile with \texttt{xelatex}.
    \item Review the multi-page PDF.
  \end{enumerate}
\end{frame}

\section{Summary}

\begin{frame}{Key Takeaways}
  Beamer turns LaTeX into a presentation engine:
  \begin{itemize}
    \item Frames are the basic unit (one per slide).
    \item Sections structure the talk.
    \item The outline frame helps the audience follow along.
  \end{itemize}
\end{frame}

\end{document}
